\PassOptionsToPackage{table,xcdraw}{xcolor}
\documentclass[11pt, a4paper]{article}

\usepackage[T1]{fontenc}
\usepackage[utf8]{inputenc}
\usepackage{tgpagella}
\usepackage[spanish, es-noshorthands]{babel}
\usepackage{amsmath, amssymb}
\usepackage{booktabs}
\usepackage{tabularx}
\usepackage{multicol}
\usepackage[most]{tcolorbox}
\usepackage[american]{circuitikz}
\usepackage[margin=2cm]{geometry}
\usepackage{fancyhdr}

\pagestyle{fancy}
\fancyhf{}
\setlength{\headheight}{16pt}
\lhead{\textbf{UCB Sede Tarija} | IMT-221}
\chead{Lab 4: Limitaciones Dinámicas (Físico)}
\rhead{Semana 10 - Sesión 1}
\lfoot{Circuitos Electrónicos II}
\rfoot{Página \thepage}
\renewcommand{\headrulewidth}{0.4pt}
\definecolor{ucbblue}{RGB}{0, 51, 102}

\begin{document}

\begin{center}
    \Large\textbf{\textcolor{ucbblue}{Guía Física Lab 4: Validación Experimental de GBW y SR}}[cite: 13]
\end{center}

\begin{tcolorbox}[colback=ucbblue!5, colframe=ucbblue, title=\textbf{Objetivo y Pregunta Experimental}]
    ¿Bajo qué combinaciones de ganancia, frecuencia y amplitud de salida un Op-Amp real se aparta de la respuesta ideal, y pueden las especificaciones de GBW y Slew Rate predecir esa desviación?[cite: 13]
\end{tcolorbox}

\section*{1. Identificación de Dispositivos y Opciones de Fuente[cite: 13]}
\begin{sloppypar}
\textbf{Advertencia Crítica:} LM741, LM358N y TL074 \textbf{NO} comparten el mismo pinout[cite: 13]. Si reemplaza un dispositivo por otro, \textbf{debe} rearmar el circuito guiado por el datasheet exacto[cite: 13].
\end{sloppypar}

\vspace{0.2cm}
\textbf{Opciones de Fuente de Alimentación[cite: 13]:}
\begin{itemize}
    \item \textbf{Opción A (Banco):} Configurar fuente dual a nominalmente $+10\,\text{V}, 0\,\text{V}, -10\,\text{V}$[cite: 13].
    \item \textbf{Opción B (Baterías):} Dos baterías de $9\,\text{V}$ en serie. \textbf{El punto medio es $0\,\text{V}$ (referencia del circuito)}[cite: 13]. Los rieles reales serán $\approx \pm 9\,\text{V}$ y pueden ser asimétricos[cite: 13].
\end{itemize}

\begin{tcolorbox}[colback=white, colframe=black, title=\textbf{Pre-Power Checklist}]
    \begin{tabularx}{\textwidth}{X X X}
    $\square$ Marcado del CI exacto & $\square$ Datasheet consultado & $\square$ Pinout verificado \\
    $\square$ Rieles medidos sin CI & $\square$ $V_+$ real = \rule{1.5cm}{0.4pt} & $\square$ $V_-$ real = \rule{1.5cm}{0.4pt} \\
    $\square$ Referencia común $0\,\text{V}$ & $\square$ Feedback verificado & $\square$ Desacoplo ($100\,\text{nF}$) puesto \\
    \end{tabularx}
\end{tcolorbox}

\section*{2. Validación Baseline de Baja Frecuencia[cite: 13]}
Antes de subir la frecuencia, debe garantizar que el amplificador funciona correctamente en DC/Baja Frecuencia[cite: 13]. Aplique una señal senoidal a $f = 1\,\text{kHz}$[cite: 13]. 
\begin{itemize}
    \item Mida la amplitud real con el osciloscopio: $V_{in,pp}$ y $V_{out,pp}$[cite: 13].
    \item Calcule la ganancia experimental: $A_{v,\mathrm{exp}} = \frac{V_{out,pp}}{V_{in,pp}}$[cite: 13].
    \item \textbf{Criterio de paso:} La onda debe ser senoidal pura, sin recortes (clipping) por los rieles de alimentación y la ganancia debe ser consistente con la topología[cite: 13].
\end{itemize}

\section*{3. Experimento A: Ancho de Banda de Pequeña Señal (GBW)[cite: 13]}
\textbf{Propósito:} Observar la caída de ganancia cerrada minimizando la demanda de Slew Rate[cite: 13].

\vspace{0.3cm}
\begin{minipage}{0.45\textwidth}
    \raggedright
    \textbf{Configuración:} Amplificador No Inversor[cite: 13].
    \begin{align*}
        R_g &= 10\,\text{k}\Omega, \quad R_f = 90\,\text{k}\Omega \\
        A_v &= 1 + \frac{90}{10} = 10 \\
        A_N &= 10 \text{ (Ganancia de Ruido)}
    \end{align*}
    \textbf{Predicción de Ancho de Banda}[cite: 13]:
    \begin{equation*}
        f_{BW} \approx \frac{GBW_{\mathrm{datasheet}}}{A_N}
    \end{equation*}
    \textbf{Condición de Entrada:} $V_{in} = 100\,\text{mV}_{pp}$ centrada en $0\,\text{V}$ (Salida esperada LF $\approx 1.0\,\text{V}_{pp}$)[cite: 13].
\end{minipage}
\hfill
\begin{minipage}{0.5\textwidth}
    \begin{center}
        \begin{circuitikz}[scale=0.8, transform shape, thick]
            \draw (0,0) node[op amp, noinv input up] (opamp) {};
            \draw (opamp.+) to[short, -o] ++(-1,0) node[left]{$V_{in}$};
            \draw (opamp.-) -- ++(0,-1.5) coordinate(fb);
            \draw (fb) to[R, l={$R_g=10\,\text{k}\Omega$}] ++(0,-2) node[ground]{};
            \draw (fb) to[R, l_={$R_f=90\,\text{k}\Omega$}] (fb -| opamp.out) -- (opamp.out);
            \draw (opamp.out) to[short, *-o] ++(1,0) node[right]{$V_{out}$};
            \draw (opamp.up) -- ++(0,0.5) node[above]{$+V$};
            \draw (opamp.down) -- ++(0,-0.5) node[below]{$-V$};
        \end{circuitikz}
    \end{center}
\end{minipage}

\vspace{0.4cm}
\textbf{Procedimiento:} Mida $V_{out,LF}$ a $1\,\text{kHz}$. Aumente $f$ hasta que $V_{out,-3dB} \approx 0.707 \cdot V_{out,LF}$. Registre datos[cite: 13].

\begin{table}[h!]
    \centering
    \footnotesize
    \renewcommand{\arraystretch}{1.3}
    \begin{tabularx}{\textwidth}{|c|c|c|c|c|c|>{\raggedright\arraybackslash}X|}
    \hline
    \textbf{Disp.} & \textbf{$f$} & \textbf{$V_{in,pp}$} & \textbf{$V_{out,pp}$} & \textbf{Ganancia ($|A_v|$)} & \textbf{Ganancia Relativa} & \textbf{Observación} \\ \hline
    & $1\,\text{kHz}$ & & & & $1.00$ & Baseline[cite: 13] \\ \hline
    & $10\,\text{kHz}$ & & & & & \\ \hline
    & $30\,\text{kHz}$ & & & & & \\ \hline
    & $100\,\text{kHz}$ & & & & & \\ \hline
    & $300\,\text{kHz}$ & & & & & \\ \hline
    & $1\,\text{MHz}$ & & & & & \\ \hline
    \end{tabularx}
\end{table}

\section*{4. Experimento B: Slew Rate de Gran Señal (SR)[cite: 13]}
\textbf{Propósito:} Aumentar la demanda de pendiente reduciendo la ganancia para evitar que el GBW limite primero[cite: 13].

\vspace{0.3cm}
\begin{minipage}{0.45\textwidth}
    \raggedright
    \textbf{Configuración:} Ganancia reducida a 2[cite: 13].
    \begin{equation*}
        R_g = 10\,\text{k}\Omega, \quad R_f = 10\,\text{k}\Omega \implies A_v = 2
    \end{equation*}
    \textbf{Derivación de la Demanda de SR}[cite: 13]: \\
    Para $v_o(t) = V_p \sin(2\pi f t)$[cite: 13]:
    \begin{equation*}
        \frac{dv_o}{dt} = 2\pi f V_p \cos(2\pi f t)
    \end{equation*}
    La pendiente máxima ocurre cuando $\cos(\cdot) = 1$[cite: 13]:
    \begin{equation*}
        \boxed{ SR_{\mathrm{req}} = 2\pi f V_p }
    \end{equation*}
\end{minipage}
\hfill
\begin{minipage}{0.5\textwidth}
    \raggedright
    \textbf{Salida Deseada a $1\,\text{kHz}$:} $V_{out} \approx 8\,\text{V}_{pp} \implies V_p = 4\,\text{V}$[cite: 13].
    
    \vspace{0.2cm}
    \textbf{Check de Saturación (Clipping)}[cite: 13]: \\
    Si a $1\,\text{kHz}$ la onda se recorta por los rieles, \textbf{debe reducir la amplitud}, registrar el nuevo $V_p$ real y recalcular todos sus $SR_{\mathrm{req}}$ de la tabla antes de seguir[cite: 13].
    
    \vspace{0.2cm}
    \textbf{Requerimientos para $V_p = 4\,\text{V}$}[cite: 13]:
    \begin{itemize}
        \item $10\,\text{kHz}: SR_{\mathrm{req}} \approx 0.251\,\text{V}/\mu\text{s}$[cite: 13]
        \item $20\,\text{kHz}: SR_{\mathrm{req}} \approx 0.503\,\text{V}/\mu\text{s}$[cite: 13]
        \item $50\,\text{kHz}: SR_{\mathrm{req}} \approx 1.257\,\text{V}/\mu\text{s}$[cite: 13]
    \end{itemize}
\end{minipage}

\vspace{0.4cm}
\begin{table}[h!]
    \centering
    \footnotesize
    \renewcommand{\arraystretch}{1.3}
    \begin{tabularx}{\textwidth}{|c|c|c|c|c|>{\centering\arraybackslash}X|>{\centering\arraybackslash}X|}
    \hline
    \textbf{Disp.} & \textbf{$f$} & \textbf{$V_{out,pp}$} & \textbf{$V_p$} & \textbf{$SR_{\mathrm{req}}$} & \textbf{Forma de Onda} & \textbf{Pendiente Medida} \\ \hline
    & $1\,\text{kHz}$ & & & & & \\ \hline
    & $5\,\text{kHz}$ & & & & & \\ \hline
    & $10\,\text{kHz}$ & & & & & \\ \hline
    & $20\,\text{kHz}$ & & & & & \\ \hline
    & $50\,\text{kHz}$ & & & & & \\ \hline
    & $100\,\text{kHz}$ & & & & & \\ \hline
    \end{tabularx}
\end{table}

\section*{5. Interpretación de Evidencia (Para el Informe IEEE)[cite: 13]}
\begin{enumerate}
    \setlength{\itemsep}{0pt}
    \item ¿Dónde aparece el punto a $-3\,\text{dB}$ y cómo se compara en orden de magnitud con $GBW/A_N$?[cite: 13]
    \item \textbf{Troubleshooting Crítico:} Antes de interpretar GBW/SR, ¿verificó que el generador y la punta del osciloscopio no introduzcan atenuación propia a altas frecuencias?[cite: 13]
    \item Diferencie: Clipping (Techo plano por rieles), GBW (Atenuación senoidal suave) y SR (Triangularización)[cite: 13].
    \item Si comparó dos dispositivos (Ej. UA741 vs TL074), ¿cuál se apartó primero y qué parámetro del datasheet justifica esto?[cite: 13]
\end{enumerate}

\end{document}
